\begin{center}
\small
\begin{longtable}{p{0.26\textwidth} p{0.31\textwidth} p{0.33\textwidth}}
\caption{Промежуточные эксперименты и выбор финальной модели}
\label{tab:intermediate-experiments}\\
\toprule
Эксперимент & Что проверялось & Итог \\
\midrule
\endfirsthead

\toprule
Эксперимент & Что проверялось & Итог \\
\midrule
\endhead

\bottomrule
\endlastfoot

Baseline CatBoost &
Компактная модель на признаках истории, региона, площади и координат поля &
\textbf{Выбран.} Зафиксирован как активная финальная модель в \path{artifacts/model_registry.json}; на test: Accuracy = 0.699006, Macro F1 = 0.568485, Weighted F1 = 0.689626 \\

CatBoost + extra features &
Расширение признакового пространства derived-признаками истории: число уникальных культур, наличие \texttt{fallow} и \texttt{legumes}, признаки повторяемости &
\textbf{Не выбран.} На test: Accuracy = 0.699235, Macro F1 = 0.569473, Weighted F1 = 0.689991. Несмотря на близкое качество, ветка сохранена как экспериментальная \texttt{improved} и не заменяет baseline в финальном проектном контракте \\

CatBoost + transition graph &
Гибридное ранжирование по формуле $\mathrm{score}=\alpha p_{\mathrm{catboost}}+\beta p_{\mathrm{graph}}$ с лучшей конфигурацией $\alpha=0.9$, $\beta=0.1$ &
\textbf{Не выбран.} Лучший hybrid-вариант на test немного уступил чистому CatBoost: Accuracy = 0.698714, Macro F1 = 0.564387, Weighted F1 = 0.688729; graph-сигнал оставлен как explainability layer \\

\end{longtable}
\footnotesize
Источники: \path{artifacts/model_registry.json},
\path{artifacts/catboost/improved/meta.json},
\path{notebooks/05_catboost_improved.ipynb},
\path{artifacts/results/classification_quality/catboost_baseline_metrics.csv},
\path{artifacts/results/graph_hybrid/hybrid_metrics.csv}.
\normalsize
\end{center}
